\section*{Résumé}

Ce rapport présente une étude menée dans le cadre du \textbf{\acrfull{PMI}}, portant sur l’évaluation de la \textbf{faisabilité
d’un système de séparation d’étages de fusée par aérofreinage}, appliqué à un \textbf{lanceur expérimental}. L’objectif du
projet n’est pas l’optimisation d’un système industrialisable, mais la \textbf{validation méthodologique et technique} d’un
principe innovant, ainsi que la mise en place des moyens nécessaires à son \textbf{évaluation expérimentale}.\\

Après une analyse de l’\textbf{état de l’art des systèmes de séparation d’étages}, le choix s’est porté sur une approche par
\textbf{aérofreinage}, reposant sur l’augmentation contrôlée de la \textbf{traînée aérodynamique} afin de créer un différentiel
de décélération entre deux étages. Ce choix a conduit à la conception d’un dispositif intégrant des \textbf{aérofreins
déployables}, ainsi qu’à la définition d’un lanceur dédié, \textbf{SP-02}, pensé comme un \textbf{vecteur expérimental} capable
de reproduire des conditions représentatives d’un lanceur multi-étages tout en supportant l’expérience étudiée.\\

La conception du système s’est appuyée sur une \textbf{maquette \acrshort{CAO} simplifiée}, utilisée à la fois pour les études
aérodynamiques et structurelles. Une \textbf{analyse par dynamique des fluides numérique (\acrshort{CFD})} a permis de
caractériser l’influence du \textbf{taux d’ouverture des aérofreins} sur le comportement aérodynamique global. Les simulations
mettent en évidence une \textbf{augmentation progressive et maîtrisée du coefficient de traînée}, associée à des phénomènes de
surpression, de \textbf{dépression}, de \textbf{décollement de la couche limite} et de \textbf{recirculation}, confirmant
l’efficacité du principe d’aérofreinage pour la séparation d’étages.\\

En parallèle, une \textbf{analyse par éléments finis (\acrshort{FEM})} a été menée afin de vérifier la \textbf{tenue mécanique}
des composants critiques soumis aux efforts aérodynamiques générés lors du déploiement des aérofreins. Cette étude a permis de
valider la \textbf{cohérence structurelle} de l’architecture retenue et d’identifier les \textbf{zones sensibles}, garantissant
la compatibilité des choix de conception avec une mise en œuvre expérimentale.\\

Un autre aspect central du projet réside dans le développement d’une \textbf{chaîne de mesure embarquée} complète, reposant sur
la plateforme \textbf{APEX}. Cette chaîne inclut l’\textbf{avionique}, la \textbf{synchronisation temporelle des capteurs}, la
\textbf{gestion des données} et les \textbf{télécommunications}, afin d’assurer l’exploitabilité des mesures lors d’un événement
fortement transitoire tel que la séparation d’étages. Elle constitue un outil essentiel pour la \textbf{validation expérimentale}
du concept et un socle réutilisable pour de futurs projets.\\

Les travaux réalisés montrent que la \textbf{séparation d’étages par aérofreinage} est techniquement envisageable dans un cadre
expérimental étudiant. Les limites identifiées, liées aux \textbf{hypothèses de modélisation}, aux conditions réelles de vol et
au \textbf{caractère transitoire du phénomène}, soulignent la nécessité des \textbf{essais en vol} pour confronter les résultats
numériques à la réalité. Ce projet constitue ainsi une \textbf{première étape structurante}, tant sur le plan technique que
méthodologique, ouvrant la voie à des développements et expérimentations futures.

\newpage

\section*{Abstract}

This report presents a study carried out as part of the \textbf{Master Innovation Project (\acrshort{PMI})}, focusing on the
assessment of the \textbf{feasibility of a rocket stage separation system based on aerodynamic braking}, applied to an
\textbf{experimental launch vehicle}. The objective of the project is not to develop an optimized or industrial solution, but to
\textbf{validate the technical and methodological relevance} of an innovative separation concept and to establish the means
required for its \textbf{experimental evaluation}.\\

Following a review of the \textbf{state of the art of stage separation systems}, an \textbf{aerodynamic braking} approach was
selected, relying on the controlled increase of \textbf{aerodynamic drag} to generate a differential deceleration between two
stages. This choice led to the design of a system incorporating \textbf{deployable airbrakes}, as well as the development of a
dedicated experimental launcher, \textbf{SP-02}, conceived as a \textbf{test vehicle} capable of reproducing representative
multi-stage flight conditions while supporting the separation experiment.\\

The system design was based on a \textbf{simplified CAD model}, used for both aerodynamic and structural analyses. A
\textbf{\acrfull{CFD}} study was conducted to characterize the influence of the \textbf{airbrake opening rate} on the overall
aerodynamic behavior. The simulations highlight a \textbf{progressive and controlled increase of the drag coefficient},
associated with \textbf{pressure gradients}, \textbf{boundary layer separation}, and \textbf{recirculation zones}, confirming
the effectiveness of aerodynamic braking for stage separation.\\

In parallel, a \textbf{\acrfull{FEM}} analysis was performed to assess the \textbf{structural integrity} of critical components
subjected to aerodynamic loads during airbrake deployment. This analysis validated the \textbf{mechanical consistency} of the
selected design and identified critical areas requiring attention for experimental implementation.\\

A major contribution of the project lies in the development of a complete \textbf{onboard measurement chain}, based on the
\textbf{APEX} platform. This system includes \textbf{avionics architecture}, \textbf{sensor synchronization}, \textbf{data
management}, and \textbf{telecommunication}, ensuring the reliability and usability of flight data during a highly transient
event such as stage separation. This instrumentation represents a key element for the \textbf{experimental validation} of the
concept and a reusable foundation for future projects.\\

The results demonstrate that \textbf{stage separation by aerodynamic braking} is \textbf{technically feasible} within a
student-level experimental framework. Identified limitations related to \textbf{modeling assumptions}, real flight conditions,
and the \textbf{transient nature of the phenomenon} emphasize the necessity of \textbf{flight testing} to confront numerical
predictions with experimental observations. Overall, this project represents a \textbf{foundational step}, both technically and
methodologically, toward the study and validation of innovative stage separation systems.
